\scanpage{61}
\uline{Corollaire} Pour que $f$ soit un iso, il f. et s. que

a) $f$ induise
\[
\pi_0(f):\ \pi_0(C)\overset{\sim}{\longrightarrow}\pi_0(C')
\]

b) $f$ induise
\[
\unsure{\operatorname{Ob}|\operatorname{O}}f=f_0:\quad\operatorname{Ob}(C)\overset{\sim}{\longrightarrow}\operatorname{Ob}(C')
\]

c) $\forall i\in I\ ;\quad f_{a_i,a_i}:\operatorname{Aut}(a_i)\overset{\sim}{\longrightarrow}\operatorname{Aut}(f_0(a_i))$

\uline{NB} $f$ \unsure{équivalence|équivalen.} $\Longleftrightarrow$ a) et c) sont vrais

\uline{Corollaire} Pour qu'une application $f_0:C_0\longrightarrow C'_0$ provienne d'un iso $f:C\simeq C'$, il f. et s. que

a) $f_0$ \deleted{soit} bijective

b) $f_0$ transforme la relation d'équivalence de $C_0$ (\unsure{pour|par} comp. \unsure{connexes}) en celle de $C'_0$ (condition vide si $C,C'$ connexes)

c) $\forall i\in I$, $\operatorname{Aut}(f_0(a_i))$ est isom à $\operatorname{Aut}(a_i)$.

\begin{center}\rule{2em}{0.3pt}\end{center}

\deleted{\unsure{Se|So}} Supposons $C=C'$. Alors on trouve
\[
\operatorname{Aut}(C)\longrightarrow\unsure{\mathfrak S|\mathscr S}_{C_0}\qquad\text{\uline{surjective}}
\]
si $\operatorname{Aut}^0(C)$ est le noyau, on trouve donc
% archive sheet number: 60
